\chapter{Implementation}
\label{ch:implementation}

\section{Introduction}
This chapter presents the implementation of the Saknny backend for the full Chapter 3 scope using the university-dorm allocation model. The delivered backend translates the functional requirements and event tables into production-ready API modules, relational data models, and operational workflows. The implementation also preserves the project governance conventions, including standard API response envelopes, role-based access control, and auditability.

\section{Implementation Architecture}
\subsection{Backend stack}
The backend is implemented using FastAPI, SQLAlchemy, and PostgreSQL. Authentication is handled through JWT bearer tokens. All APIs are exposed under \texttt{/api/v1} and follow a unified response format:
\begin{verbatim}
{
  "success": true|false,
  "data": {...} | null,
  "error": null | "error message"
}
\end{verbatim}

\subsection{Module organization}
The implementation is divided into functional endpoint modules:
\begin{itemize}
  \item Identity and verification: \texttt{auth.py}, \texttt{students.py}, \texttt{admin.py}
  \item Catalog: \texttt{catalog.py}
  \item Applications and waitlist: \texttt{applications.py}
  \item Allocation: \texttt{allocations.py}
  \item Contracts and lease lifecycle: \texttt{leases.py}
  \item Billing simulation: \texttt{billing.py}
  \item Check-in and residency lifecycle: \texttt{residency.py}
  \item Maintenance workflows: \texttt{maintenance.py}
  \item Messaging and announcements: \texttt{communications.py}
  \item Analytics and audit reporting: \texttt{insights.py}
  \item Surveys: \texttt{surveys.py}
\end{itemize}

\section{Database Implementation}
\subsection{Core schema}
The schema was extended from authentication and verification tables to a full operational model including:
\begin{itemize}
  \item \texttt{buildings}, \texttt{rooms}
  \item \texttt{applications}, \texttt{application\_reviews}
  \item \texttt{allocations}, \texttt{leases}
  \item \texttt{payment\_intents}
  \item \texttt{checkins}, \texttt{room\_change\_requests}
  \item \texttt{maintenance\_tickets}
  \item \texttt{announcements}, \texttt{messages}
  \item \texttt{audit\_logs}
  \item \texttt{surveys}, \texttt{survey\_dispatches}
\end{itemize}

\subsection{Constraints and business rules}
The implementation enforces key Chapter 3 rules at model and API layers:
\begin{itemize}
  \item Gender compatibility between student and building for allocations.
  \item Room capacity constraints (\texttt{available\_beds} range and decrement on assignment).
  \item Enrollment gate before application submission.
  \item One active allocation per student.
  \item Lease issued only from valid assigned allocation.
  \item Controlled enum states for verification, applications, leases, payments, tickets, and lifecycle requests.
\end{itemize}

\subsection{Migration workflow}
Alembic was added to version schema evolution:
\begin{itemize}
  \item Configuration: \texttt{backend/alembic.ini}
  \item Environment: \texttt{backend/alembic/env.py}
  \item Initial full migration: \texttt{backend/alembic/versions/20260515\_01\_full\_ch3\_schema.py}
\end{itemize}
This replaces the previous non-versioned schema approach as the primary path for database initialization.

\section{API Implementation by Functional Area}
\subsection{Identity and verification}
Students can register, update profiles, upload verification documents, and view document status. Admins can review documents and update enrollment status with audit logging.

\subsection{Catalog and applications}
Catalog APIs provide filtered building/room browsing and admin inventory management. Application APIs implement submission, review, finalization, and waitlist workflows.

\subsection{Allocation, contracts, billing}
Approved applications can be allocated to available rooms with rule checks. Leases are issued and signed digitally. Payments are modeled through a simulated gateway lifecycle (\texttt{initiated} to \texttt{paid/failed/refunded}).

\subsection{Residency lifecycle and maintenance}
The residency module supports check-in, key issuance, room-change requests, and checkout. Maintenance workflows include ticket creation, assignment, escalation, and resolution tracking.

\subsection{Communication, analytics, and surveys}
In-app messaging and targeted announcements were implemented for stakeholder communication. Admin dashboards expose occupancy, payment, and ticket metrics. Survey creation, dispatch, and completion are supported for satisfaction indicators.

\section{Auditability and security}
All sensitive state transitions write audit records through a shared audit service. RBAC is enforced through dependency guards for student and admin actions. Ownership checks are implemented for student-scoped resources (documents, leases, payments, surveys).

\section{Requirement Mapping}
A complete requirement-to-implementation traceability matrix is provided in:
\begin{itemize}
  \item \texttt{docs/functional\_traceability\_matrix.md}
\end{itemize}
This matrix maps Chapter 3 event groups to endpoints, data models, and implemented validation logic.

\section{Testing and verification}
Automated coverage checks were added in:
\begin{itemize}
  \item \texttt{backend/tests/test\_traceability.py}
\end{itemize}
The tests verify:
\begin{itemize}
  \item Presence of all core Chapter 3 API paths in router registration.
  \item Compliance of response envelope helpers.
  \item Presence of required Chapter 3 tables in the database schema contract.
\end{itemize}

\section{ML readiness and extension points}
By request, no ML model was implemented in this phase. However, the backend is prepared for future ML integration through:
\begin{itemize}
  \item Student preference and distance-related fields.
  \item Structured historical data for applications, allocations, payments, and maintenance.
  \item Auditable event logs suitable for feature engineering and model evaluation.
\end{itemize}
Future work can add recommendation, prioritization, or anomaly detection services without schema redesign.

\section{Conclusion}
The delivered implementation operationalizes the full Chapter 3 backend scope for the university-dorm model, with complete API modules, expanded relational schema, migration support, auditability, and documentation artifacts for report integration.
